% GENERATED FILE. Edit canonical lessons, then run npm run generate:books.
% canonical-source-hash: fnv1a64:75c69021a077b14e
% canonical-lessons: PT-C29-palavras, PT-C29-linhas, PT-C29-primeira-leitura

\chapter{Reading — Words, Lines, and a First Passage}
\label{ch:pt-primeira-leitura}
\hlchaptermodality{29}

\begin{chapteropening}
\textbf{By the end of this chapter:} \emph{I can read a Portuguese word without its article, follow a line whose subject lives inside the verb, and hold a whole passage without translating it word by word.}

Forty-two words of Portuguese held together by one joining word and the order of its sentences, because the track has not yet taught but or because.
\end{chapteropening}

\section[(reading)]{(reading) — six words off a counter}

\label{lesson:PT-C29-palavras}



\begin{quote}
Every word below you have met. Not one is new.

What is new is that nobody says them first.
\end{quote}

\begin{tcolorbox}[breakable,skin=enhanced,colback=orange!6,colframe=orange!45!black,boxrule=0.5pt,arc=1mm,left=6pt,right=6pt,top=4pt,bottom=4pt,fonttitle=\bfseries,title={Read this}]
\begin{quote}
água pão café leite livro amigo
\end{quote}

Read down the list once, without stopping.

Again — and this time say the article you would put in front of each one. \textbf{A} água. \textbf{O} pão.
\end{tcolorbox}

\begin{culture}
Portuguese hangs an article on nearly every noun, and the article is the gender said out loud.

That is why a written list like this one is harder than a real Portuguese page. On a page the article is already there, doing the remembering for you. Stripped of it, each word makes you do the work yourself.

Six words, and the hardest thing about them is what is missing.
\end{culture}

\subsection*{Your turn}
\begin{itemize}
\raggedright
  \item \emph{Say it:} the six words with their articles, once, without stopping
\end{itemize}

\emph{Twice through:}

\begin{tcolorbox}[breakable,colback=teal!4,colframe=teal!35!black,title={Before you move on}]
How many of the six were new? (\textbf{None.})

What is missing from this list that a real page would have? (\textbf{The article} — and it carries the gender.)
\end{tcolorbox}
\section[(reading)]{(reading) — six lines, and one that depends on you}

\label{lesson:PT-C29-linhas}



\begin{quote}
The last lesson gave you words. These are things you would say.

You have said every one of them.
\end{quote}

\begin{tcolorbox}[breakable,skin=enhanced,colback=orange!6,colframe=orange!45!black,boxrule=0.5pt,arc=1mm,left=6pt,right=6pt,top=4pt,bottom=4pt,fonttitle=\bfseries,title={Read this}]
\begin{quote}
Olá! Bom dia. Boa noite. Obrigado. Sim. Adeus!
\end{quote}

Read down once, without stopping.

Again — and notice that \textbf{Bom dia} and \textbf{Boa noite} differ by one letter, and that the letter is doing grammar.
\end{tcolorbox}

\begin{culture}
\textbf{Bom} and \textbf{Boa} are the same word agreeing with \textbf{dia} and \textbf{noite}. Read the two lines together and the agreement is visible without anybody naming it.

\textbf{Obrigado} is the strange one. It means \emph{obliged}, and it is a participle, so it agrees with the person SAYING it, not the person thanked. A woman writes \textbf{obrigada}.

That is the only line on this page whose spelling depends on who is reading it aloud.
\end{culture}

\subsection*{Your turn}
\begin{itemize}
\raggedright
  \item \emph{Say it:} the six lines, once, in order, without stopping
\end{itemize}

\emph{Twice through:}

\begin{tcolorbox}[breakable,colback=teal!4,colframe=teal!35!black,title={Before you move on}]
Why is it \textbf{Bom dia} but \textbf{Boa noite}? (\textbf{Agreement} — and you read it before it was explained.)

Who does \textbf{obrigado} agree with? (\textbf{The speaker.})
\end{tcolorbox}
\section[(reading)]{(reading) — the first passage}

\label{lesson:PT-C29-primeira-leitura}



\begin{quote}
Six words, then six lines. Now twelve of them together.

Forty-two words, and not one of them is new.
\end{quote}

\begin{tcolorbox}[breakable,skin=enhanced,colback=orange!6,colframe=orange!45!black,boxrule=0.5pt,arc=1mm,left=6pt,right=6pt,top=4pt,bottom=4pt,fonttitle=\bfseries,title={Read this}]
\begin{quote}
Olá! Bom dia. Eu tenho um livro. O livro é muito bom. Eu leio e falo. Tenho um amigo. O amigo é o primeiro. A família é boa. Hoje eu como pão. Tenho água e café. Amanhã leio o livro. Obrigado. Adeus!
\end{quote}

Read it once without stopping. Do not translate as you go — let the sentences arrive.

Now read it again, and count the joining words.
\end{tcolorbox}

\begin{culture}
One. \textbf{E}, twice.

There is no \emph{but} and no \emph{because} here, and that is not a simplification — the book has not taught either yet. What holds the paragraph together instead is the order of the sentences: a thing, then what it is like; today, then tomorrow.

That is worth knowing before the joiners arrive. A paragraph can be held together by sequence alone, and every language lets you do it.

Notice also that \textbf{eu} appears and then stops. Portuguese keeps the person in the verb, so \textbf{Tenho} on its own is already \emph{I have}, and the paragraph drops the pronoun once it is established.

This is the first whole passage the book has asked you to hold at once, and it works because everything in it was paid for.
\end{culture}

\subsection*{Your turn}
\begin{itemize}
\raggedright
  \item \emph{Say it:} the passage aloud, once, without stopping
\end{itemize}

\emph{Twice through:}

\begin{tcolorbox}[breakable,colback=teal!4,colframe=teal!35!black,title={Before you move on}]
What holds this paragraph together, if not joining words? (\textbf{The order of the sentences.})

Why does \textbf{eu} disappear halfway? (\textbf{The verb already carries it.})
\end{tcolorbox}
